
sur $\underline{k}(e)$, car ces
points sont en correspondance biunivoque avec les drapeaux du type $(1,\ldots,1)$ ($n$ termes) de l'espace vectoriel $\underline{E}_e$ de dimension $n>0$ sur $\underline{k}(e)$. D'autre part ces points rationnels correspondent aussi aux germes de sections de $g$ au voisinage de $e$, de sorte qu'on peut trouver une telle section $s:U\to D$ au-dessus d'un ouvert $U$ de $X$. Soit $X'=\overline{s(U)}$ l'adhérence de $s(U)$ munie de la structure induite réduite, et soit $f:X'\to X$ la restriction de $g$ à $X'$. Le morphisme $f$ est projectif, et comme il est un isomorphisme sur $s(U)$, il est birationnel.
Donc $\dim(X')\leq\dim(X)$ (EGA IV 5.6.6.1). D'autre part $f(X')=g(X')$ est fermé et contient $U$, de sorte que $f$ est surjectif et $\dim(X')\geq\dim(X)$ (EGA IV 5.4.1(ii)).

Les hypothèses du cas I sont donc satisfaites pour $X'$ et $x'=f^*(x)\in\Filt^i(X')$, et on trouve
\begin{equation}\tag{1.3.7'}
f^*(x)\cl_\bullet(\underline{O}_{X'})\in\Filt_{j-i}(X').
\end{equation}
On en conclut, tenant compte de la formule de projection et de (1.1.8) :
\begin{equation}\tag{1.3.8}
x f_*(\cl_\bullet(\underline{O}_{X'}))\in\Filt_{j-i}(X).
\end{equation}

D'autre part l'homomorphisme canonique
\[
\underline{O}_X\longrightarrow f_*(\underline{O}_{X'})
\]
est un isomorphisme sur l'ouvert dense $U$, et les $\underline{O}_X$-modules $R^if_*(\underline{O}_{X'})$ ($i\geq1$) sont concentrés sur le sous-ensemble fermé rare $X-U$. Il s'ensuit que
\begin{equation}\tag{1.3.9}
f_*(\cl_\bullet(\underline{O}_{X'}))=\cl_\bullet(\underline{O}_X)+w,
